\documentclass[11pt]{article}
\usepackage[margin=1in]{geometry}
\usepackage{booktabs,amsmath,hyperref}
\title{Scope-Correct Global Filtering for a Fleet Trace Investigation Interface}
\author{AssistX fleet observability research draft}
\date{October 8, 2026}
\begin{document}
\maketitle
\begin{abstract}
Paginated observability dashboards may misstate incident prevalence when outcome filtering is performed only on an in-memory page of 50 traces. This study proposes and tests an authenticated, read-only global outcome-filtering extension for an indexed fleet trace graph. The implementation applies the same existence predicates to both a filtered count and a paginated query, gives failed events precedence over completed events, rejects unsupported filters, and bounds query execution time. Offline synthetic and contract tests passed (17 Python and 11 JavaScript). Live 85,000-group query performance and browser accessibility remain unverified.
\end{abstract}
\section{Problem and prediction}
The previous workbench showed an all-history correlation-ID count alongside current-page status counts. An outcome selector limited to the loaded page could not find failures elsewhere in the database. Before implementation we preregistered a global failed/completed/open filter, identical membership rules in count and page queries, parameterized search, a four-second query budget, filtered permalinks, and explicit failure when connected to an older server.
\section{Predicate method}
For a trace group $g$ with event-type set $T(g)$, define $F(g)$ as existence of an event type ending in \texttt{.failed}, and $C(g)$ as existence of one ending in \texttt{.completed} or \texttt{.accepted}. On nonempty event groups we use:
\[
O(g) = \begin{cases}
\mathrm{failed}, & F(g),\\
\mathrm{completed}, & \neg F(g)\wedge C(g),\\
\mathrm{open}, & \neg F(g)\wedge\neg C(g).
\end{cases}
\]
The implementation uses fixed Cypher \texttt{EXISTS} subqueries for all-history filtering and binds correlation-ID search parameters. The Neo4j manual documents existential subqueries \cite{neo4jexists}. The endpoint projects no payload, arguments or outputs. The filtered global total and page results share the same predicates; the UI still labels the separate current-page failure count honestly.
\section{Synthetic verification}
\begin{table}[h]\centering
\caption{Offline implementation checks; not live-fleet performance acceptance.}
\begin{tabular}{ll}\toprule
Measure & Result\\\midrule
Python query and FastAPI route tests & 17 passed\\
JavaScript workbench contract tests & 11 passed\\
Invalid filters and oversized searches & Rejected\\
Older backend lacking filter acknowledgment & Unavailable, not false zero\\
Live graph query latency and concurrency & Not measured\\
Authenticated browser and WCAG audit & Not measured\\
\bottomrule\end{tabular}\end{table}
\section{Limitations and next gates}
A query timeout does not prove the global count is cheap enough for production traffic. Release requires an operator-approved staging query plan, index inspection, large fixture latency, rate-limit review, and real keyboard/mobile browser testing. WCAG 2.2 describes focus and navigation success criteria \cite{wcag22}. Trace-group outcome metadata must not be confused with full-fidelity tool-call retention, whose preservation obligations are distinct \cite{nist92}. This manuscript is an arXiv-style draft, not compiled, submitted or peer reviewed.
\begin{thebibliography}{9}
\bibitem{neo4jexists} Neo4j, \emph{EXISTS subqueries}, Cypher Manual. \url{https://neo4j.com/docs/cypher-manual/current/subqueries/existential/}.
\bibitem{wcag22} W3C, \emph{Web Content Accessibility Guidelines 2.2}. \url{https://www.w3.org/TR/WCAG22/}.
\bibitem{nist92} K. Kent and M. Souppaya, \emph{Guide to Computer Security Log Management}, NIST SP 800-92 (2006). \url{https://doi.org/10.6028/NIST.SP.800-92}.
\end{thebibliography}
\end{document}
